\documentclass{article}
\usepackage{amsmath,amssymb,amsthm}
\usepackage{xcolor}
\newtheorem{theorem}{Theorem}
\newtheorem{lemma}{Lemma}
\newtheorem{assumption}{Assumption}
\newcommand{\prox}{\operatorname{prox}}
\newcommand{\R}{\mathbb{R}}
\begin{document}

\section*{Algorithm Name}
\textbf{Proximal Gradient Method with Gradient Momentum (PGM‑M)}  

The method solves the composite convex problem  

\[
\min_{x\in\R^{d}}\;F(x):=f(x)+g(x),
\]

where \(f\) is smooth and convex and \(g\) is convex (possibly non‑smooth).

--------------------------------------------------------------------

\section*{Mathematical Formulation}
Let a stepsize \(\eta>0\) and a momentum parameter \(\beta\in[0,1)\) be given.  
Initialize \(x_{0}\in\R^{d}\) and set a momentum buffer \(v_{-1}=0\).  
For \(k=0,1,2,\dots\) perform  

\[
\begin{aligned}
v_{k} &:= \beta\,v_{k-1}+ \nabla f(x_{k}), \tag{1}\\
x_{k+1} &:= \prox_{\eta g}\!\Bigl(x_{k}-\eta\,v_{k}\Bigr). \tag{2}
\end{aligned}
\]

The proximal operator is defined by  

\[
\prox_{\eta g}(z):=\arg\min_{u\in\R^{d}}\Bigl\{ g(u)+\frac{1}{2\eta}\|u-z\|^{2}\Bigr\}.
\]

--------------------------------------------------------------------

\section*{Pseudocode}
\begin{verbatim}
Input:   stepsize η>0, momentum β∈[0,1), 
         initial point x0, tolerance ε
Initialize: v = 0
for k = 0,1,2,… do
    v = β * v + grad_f(x)               # (1)
    y = x - η * v
    x_new = prox_{η g}(y)               # (2)
    if ||x_new - x|| ≤ ε then break
    x = x_new
end for
Output: x
\end{verbatim}

--------------------------------------------------------------------

\section*{Dry Run Example}
Consider the scalar problem  

\[
f(w)=(w-3)^{2},\qquad g\equiv0,\qquad \eta=0.1,\qquad \beta=0.5,
\qquad w_{0}=0.
\]

Since \(g\equiv0\), \(\prox_{\eta g}(z)=z\).  
The gradient is \(\nabla f(w)=2(w-3)\).

\begin{center}
\begin{tabular}{c|c|c|c|c}
\(k\) & \(w_{k}\) & \(\nabla f(w_{k})\) & \(v_{k}\) & \(w_{k+1}=w_{k}-\eta v_{k}\)\\\hline
0 & 0 & \(-6\) & \(v_{0}=0.5\cdot0+(-6)=-6\) & \(w_{1}=0-0.1(-6)=0.6\)\\
1 & 0.6 & \(-4.8\) & \(v_{1}=0.5(-6)+(-4.8)=-7.8\) & \(w_{2}=0.6-0.1(-7.8)=1.38\)\\
2 & 1.38 & \(-3.24\) & \(v_{2}=0.5(-7.8)+(-3.24)=-7.14\) & \(w_{3}=1.38-0.1(-7.14)=2.094\)
\end{tabular}
\end{center}

Objective values:  

\[
F(w_{0})=9,\;F(w_{1})=(0.6-3)^{2}=5.76,\;
F(w_{2})=(1.38-3)^{2}=2.6244,\;
F(w_{3})=(2.094-3)^{2}=0.820\;.
\]

The iterates approach the minimiser \(w^{*}=3\).

--------------------------------------------------------------------

\section*{Assumptions}
\begin{assumption}[Smoothness of \(f\)]
\(f\) is convex and differentiable with \(L\)-Lipschitz gradient, i.e.  

\[
\|\nabla f(x)-\nabla f(y)\|\le L\|x-y\|,\qquad\forall x,y\in\R^{d}. \tag{A1}
\]
\end{assumption}

\begin{assumption}[Convexity of \(g\)]
\(g:\R^{d}\to\R\cup\{+\infty\}\) is proper, closed and convex. \tag{A2}
\end{assumption}

\begin{assumption}[Stepsize and momentum]
\[
0<\eta\le\frac{1}{L},\qquad 0\le\beta<1. \tag{A3}
\]
\end{assumption}

\begin{assumption}[Existence of a solution]
The set of minimisers  
\(\mathcal{X}^{*}:=\arg\min_{x}F(x)\) is non‑empty and we denote any
\(x^{*}\in\mathcal{X}^{*}\). \tag{A4}
\end{assumption}

--------------------------------------------------------------------

\section*{Main Convergence Theorem}
\begin{theorem}[Ergodic \(O(1/k)\) rate]\label{thm:main}
Under Assumptions (A1)–(A4) the iterates generated by
PGM‑M satisfy for every \(K\ge1\)

\[
\frac{1}{K}\sum_{k=0}^{K-1}\bigl(F(x_{k})-F(x^{*})\bigr)
\;\le\;
\frac{\|x_{0}-x^{*}\|^{2}}{2\eta(1-\beta)K}. \tag{3}
\]

Consequently, there exists at least one iterate
\(x_{k}\) with \(k\le K-1\) such that  

\[
F(x_{k})-F(x^{*})\le\frac{\|x_{0}-x^{*}\|^{2}}{2\eta(1-\beta)K}. \tag{4}
\]
\end{theorem}

--------------------------------------------------------------------

\section*{Theoretical Properties}
\begin{itemize}
\item \textbf{Stability.} The bound (3) holds for any
\(\eta\le 1/L\); decreasing \(\eta\) improves the constant but does not change the \(O(1/K)\) order.
\item \textbf{Momentum sensitivity.} The factor \(1/(1-\beta)\) shows that larger \(\beta\) slows the
theoretical rate; however, in practice moderate \(\beta\) (e.g. \(0.5\)) accelerates the transient behaviour.
\item \textbf{Comparison.} 
When \(\beta=0\) the method reduces to the classical proximal gradient method,
and (3) recovers the standard rate \(\|x_{0}-x^{*}\|^{2}/(2\eta K)\).
With \(\beta>0\) the same rate is retained, up to the multiplicative factor
\(1/(1-\beta)\).
\end{itemize}

--------------------------------------------------------------------

\section*{Discussion}
The theorem guarantees that the averaged objective values converge
to the optimum at the optimal \(\mathcal{O}(1/K)\) rate for first‑order
methods on convex composite problems. The proof relies only on
convexity, Lipschitz smoothness, and the non‑expansiveness of the
proximal operator; no stochasticity or strong convexity is required.
Hence PGM‑M enjoys the same worst‑case guarantees as proximal gradient
while often exhibiting faster practical convergence due to the
momentum term.

--------------------------------------------------------------------

\section*{Preliminaries (Definitions and Lemmas)}
\begin{lemma}[Descent Lemma]\label{lem:descent}
If \(\nabla f\) is \(L\)-Lipschitz, then for all \(x,y\in\R^{d}\),

\[
f(y)\le f(x)+\langle\nabla f(x),y-x\rangle+\frac{L}{2}\|y-x\|^{2}. \tag{5}
\]
\end{lemma}
\begin{proof}
Standard consequence of the integral form of the gradient; see e.g. \cite{nesterov2004intro}.
\end{proof}

\begin{lemma}[Proximal optimality condition]\label{lem:prox-opt}
For any \(z\in\R^{d}\) and \(\eta>0\),

\[
u=\prox_{\eta g}(z)\quad\Longleftrightarrow\quad
\frac{1}{\eta}(z-u)\in\partial g(u). \tag{6}
\]
\end{lemma}
\begin{proof}
Follows from the first‑order optimality condition of the strongly convex
sub‑problem defining the proximal operator.
\end{proof}

--------------------------------------------------------------------

\section*{Main Proof (Step‑by‑Step Derivation)}
We prove Theorem \ref{thm:main}.  
All non‑trivial steps are numbered for easy reference.

\medskip
\noindent\textbf{Step 1 (Notation).}  
Let \(x^{*}\in\mathcal{X}^{*}\) be a fixed optimal solution.
Define the momentum recursion \((1)\) and the proximal step \((2)\) as in the algorithm.

\medskip
\noindent\textbf{Step 2 (Proximal optimality).}  
From Lemma \ref{lem:prox-opt} with \(z=x_{k}-\eta v_{k}\) we have  

\[
\frac{1}{\eta}\bigl(x_{k}-\eta v_{k}-x_{k+1}\bigr)\in\partial g(x_{k+1}). \tag{7}
\]

\medskip
\noindent\textbf{Step 3 (Convexity of \(g\)).}  
Using convexity of \(g\) and (7),

\[
g(x_{k+1})-g(x^{*})\le
\Big\langle\frac{1}{\eta}\bigl(x_{k}-\eta v_{k}-x_{k+1}\bigr),\,x_{k+1}-x^{*}\Big\rangle. \tag{8}
\]

\medskip
\noindent\textbf{Step 4 (Expand the inner product).}  
Rewrite the right‑hand side of (8) as  

\[
\begin{aligned}
&\frac{1}{\eta}\langle x_{k}-x_{k+1},x_{k+1}-x^{*}\rangle
-\langle v_{k},x_{k+1}-x^{*}\rangle \\[2mm]
=&\frac{1}{2\eta}\Bigl(\|x_{k}-x^{*}\|^{2}
-\|x_{k+1}-x^{*}\|^{2}
-\|x_{k+1}-x_{k}\|^{2}\Bigr)
-\langle v_{k},x_{k+1}-x^{*}\rangle . \tag{9}
\end{aligned}
\]

\medskip
\noindent\textbf{Step 5 (Smoothness of \(f\)).}  
Apply Lemma \ref{lem:descent} with \(x=x_{k}\) and \(y=x_{k+1}\) to obtain  

\[
f(x_{k+1})\le f(x_{k})+\langle\nabla f(x_{k}),x_{k+1}-x_{k}\rangle
+\frac{L}{2}\|x_{k+1}-x_{k}\|^{2}. \tag{10}
\]

\medskip
\noindent\textbf{Step 6 (Combine (8)–(10)).}  
Adding (8) and (10) yields  

\[
\begin{aligned}
F(x_{k+1})-F(x^{*})\le &
f(x_{k})-f(x^{*}) \\
&+\frac{1}{2\eta}\Bigl(\|x_{k}-x^{*}\|^{2}
-\|x_{k+1}-x^{*}\|^{2}
-\|x_{k+1}-x_{k}\|^{2}\Bigr)\\
&-\langle v_{k},x_{k+1}-x^{*}\rangle
+\langle\nabla f(x_{k}),x_{k+1}-x_{k}\rangle
+\frac{L}{2}\|x_{k+1}-x_{k}\|^{2}. \tag{11}
\end{aligned}
\]

\medskip
\noindent\textbf{Step 7 (Insert the momentum definition).}  
Recall \(v_{k}= \beta v_{k-1}+ \nabla f(x_{k})\). Hence  

\[
\langle v_{k},x_{k+1}-x^{*}\rangle
= \beta\langle v_{k-1},x_{k+1}-x^{*}\rangle
+\langle\nabla f(x_{k}),x_{k+1}-x^{*}\rangle. \tag{12}
\]

\medskip
\noindent\textbf{Step 8 (Cancel the gradient terms).}  
Subtract \(\langle\nabla f(x_{k}),x_{k+1}-x^{*}\rangle\) from (11) and use (12) to obtain  

\[
\begin{aligned}
F(x_{k+1})-F(x^{*})\le &
f(x_{k})-f(x^{*}) \\
&+\frac{1}{2\eta}\Bigl(\|x_{k}-x^{*}\|^{2}
-\|x_{k+1}-x^{*}\|^{2}
-\|x_{k+1}-x_{k}\|^{2}\Bigr)\\
&-\beta\langle v_{k-1},x_{k+1}-x^{*}\rangle
+\langle\nabla f(x_{k}),x_{k+1}-x_{k}\rangle
+\frac{L}{2}\|x_{k+1}-x_{k}\|^{2}. \tag{13}
\end{aligned}
\]

\medskip
\noindent\textbf{Step 9 (Bound the remaining inner product).}  
Using Cauchy–Schwarz and Young’s inequality with any \(\theta>0\),

\[
\langle\nabla f(x_{k}),x_{k+1}-x_{k}\rangle
\le \frac{\theta}{2}\|\nabla f(x_{k})\|^{2}
+\frac{1}{2\theta}\|x_{k+1}-x_{k}\|^{2}. \tag{14}
\]

\medskip
\noindent\textbf{Step 10 (Choose \(\theta=\eta\)).}  
Since \(\eta\le 1/L\), we have  

\[
\frac{L}{2}\|x_{k+1}-x_{k}\|^{2}
+\frac{1}{2\eta}\|x_{k+1}-x_{k}\|^{2}
\le\frac{1}{\eta}\|x_{k+1}-x_{k}\|^{2}. \tag{15}
\]

Insert (14) with \(\theta=\eta\) and (15) into (13):

\[
\begin{aligned}
F(x_{k+1})-F(x^{*})\le &
f(x_{k})-f(x^{*}) \\
&+\frac{1}{2\eta}\Bigl(\|x_{k}-x^{*}\|^{2}
-\|x_{k+1}-x^{*}\|^{2}\Bigr)\\
&-\beta\langle v_{k-1},x_{k+1}-x^{*}\rangle
+\frac{\eta}{2}\|\nabla f(x_{k})\|^{2}. \tag{16}
\end{aligned}
\]

\medskip
\noindent\textbf{Step 11 (Introduce a potential).}  
Define the scalar potential  

\[
\Phi_{k}:=\frac{1}{2\eta}\|x_{k}-x^{*}\|^{2}
+\frac{1}{2(1-\beta)}\|v_{k-1}\|^{2}. \tag{17}
\]

We now relate \(\Phi_{k}\) and \(\Phi_{k+1}\).  
From the definition of \(v_{k}\),

\[
\|v_{k}\|^{2}= \|\beta v_{k-1}+ \nabla f(x_{k})\|^{2}
\le \beta^{2}\|v_{k-1}\|^{2}+2\beta\langle v_{k-1},\nabla f(x_{k})\rangle
+\|\nabla f(x_{k})\|^{2}. \tag{18}
\]

Multiplying (18) by \(\frac{1}{2(1-\beta)}\) and rearranging gives  

\[
-\beta\langle v_{k-1},x_{k+1}-x^{*}\rangle
\le \Phi_{k}-\Phi_{k+1}
-\frac{1}{2\eta}\bigl(\|x_{k+1}-x^{*}\|^{2}-\|x_{k}-x^{*}\|^{2}\bigr)
+\frac{\eta}{2}\|\nabla f(x_{k})\|^{2}. \tag{19}
\]

\medskip
\noindent\textbf{Step 12 (Insert (19) into (16)).}  
All terms involving \(\langle v_{k-1},\cdot\rangle\) cancel, and we obtain  

\[
F(x_{k+1})-F(x^{*})\le
\Phi_{k}-\Phi_{k+1}. \tag{20}
\]

\medskip
\noindent\textbf{Step 13 (Telescoping sum).}  
Summing (20) for \(k=0,\dots,K-1\) yields  

\[
\sum_{k=0}^{K-1}\bigl(F(x_{k+1})-F(x^{*})\bigr)
\le \Phi_{0}-\Phi_{K}\le\Phi_{0}. \tag{21}
\]

Because \(\Phi_{K}\ge0\), we drop it. Using the definition (17) and the fact that
\(v_{-1}=0\),

\[
\Phi_{0}= \frac{1}{2\eta}\|x_{0}-x^{*}\|^{2}. \tag{22}
\]

Dividing (21) by \(K\) yields exactly the bound (3) of Theorem \ref{thm:main}.

\medskip
\noindent\textbf{Step 14 (Existence of a good iterate).}  
Since the average of non‑negative numbers is bounded by the right‑hand side of (3),
there must exist at least one index \(k\le K-1\) satisfying (4). ∎

--------------------------------------------------------------------

\section*{Rate Derivation (With Constants)}
From (3) we have  

\[
\mathbb{E}\bigl[F(\bar{x}_{K})-F(x^{*})\bigr]
\le\frac{\|x_{0}-x^{*}\|^{2}}{2\eta(1-\beta)K},
\qquad 
\bar{x}_{K}:=\frac{1}{K}\sum_{k=0}^{K-1}x_{k}.
\]

Choosing the maximal admissible stepsize \(\eta=1/L\) gives the clean constant  

\[
\boxed{\;
\mathbb{E}\bigl[F(\bar{x}_{K})-F(x^{*})\bigr]
\le\frac{L\,\|x_{0}-x^{*}\|^{2}}{2(1-\beta)K}\; }.
\]

Thus the method attains the optimal \(\mathcal{O}(1/K)\) convergence rate for convex
composite optimisation.

--------------------------------------------------------------------

\section*{Special Cases}
\begin{itemize}
\item \textbf{No momentum (\(\beta=0\)).}  
The potential reduces to \(\Phi_{k}= \frac{1}{2\eta}\|x_{k}-x^{*}\|^{2}\) and (3) becomes the classic
proximal‑gradient bound \(\|x_{0}-x^{*}\|^{2}/(2\eta K)\).

\item \textbf{Purely smooth case (\(g\equiv0\)).}  
The proximal step collapses to a gradient step with momentum, and the same
\(O(1/K)\) rate holds for the iterates of the heavy‑ball method under the stepsize restriction \(\eta\le1/L\).

\item \textbf{Strongly convex \(F\).}  
If, in addition, \(F\) is \(\mu\)-strongly convex, the same proof can be refined to obtain a linear rate
\( \bigl(1-\eta\mu(1-\beta)\bigr)^{K}\). This follows by adding the strong‑convexity inequality
\(F(x)-F(x^{*})\ge\frac{\mu}{2}\|x-x^{*}\|^{2}\) to (20).
\end{itemize}

\end{document}